% Vacuum Pressure Impregnation (VPI) process: six sequential steps.
\documentclass{article}
\usepackage[paperwidth=120cm,paperheight=120cm,margin=10mm]{geometry}
\pagestyle{empty}
\usepackage{fontspec}
\usepackage[bidi=basic]{babel}
\babelprovide[import]{english}
\babelprovide[main, import]{persian}
\babelfont{rm}[Renderer=Node, Path=../fonts/, Extension=.ttf, UprightFont=*-Regular, BoldFont=*-Bold]{Vazirmatn}
\babelfont{sf}[Renderer=Node, Path=../fonts/, Extension=.ttf, UprightFont=*-Regular, BoldFont=*-Bold]{Vazirmatn}
\providecommand{\setRTL}{}\providecommand{\setLTR}{}
\providecommand{\RL}[1]{\foreignlanguage{persian}{#1}}
\providecommand{\LR}[1]{\babelsublr{#1}}
\usepackage{tikz}
\usepackage{pgfplots}
\pgfplotsset{compat=1.18}
\usetikzlibrary{arrows.meta, positioning, shapes.geometric, shapes.misc, calc, fit, backgrounds, decorations.pathmorphing, patterns, angles, quotes}
\begin{document}
\setRTL
\tikzset{every node/.append style={execute at begin node=\setRTL}}
\babelsublr{\begin{tikzpicture}[>=Stealth, font=\small,
    vstep/.style={rectangle, rounded corners=3pt, draw=teal!60!black, very thick,
      fill=teal!8, text width=6.6cm, minimum height=1.0cm, align=left, inner sep=5pt},
    num/.style={circle, draw=teal!60!black, fill=teal!20, very thick,
      minimum size=0.7cm, inner sep=0pt, font=\bfseries}]

  \node[font=\bfseries] (title) at (0,1.0) {مراحل فرایند VPI};

  \def\dy{1.35}
  \node[vstep] (a) at (0,0)        {قرار دادن سیم‌پیچ در محفظهٔ خلأ};
  \node[vstep] (b) at (0,-1*\dy)   {تخلیهٔ هوا و رطوبت (خلأ)، چند ساعت};
  \node[vstep] (c) at (0,-2*\dy)   {تزریق رزین اپوکسی کم‌گرانروی زیر خلأ};
  \node[vstep] (d) at (0,-3*\dy)   {اعمال فشار گاز برای راندن رزین به حفره‌ها};
  \node[vstep] (e) at (0,-4*\dy)   {پخت گرمایی (نیم‌رخ دما / زمان)};
  \node[vstep] (f) at (0,-5*\dy)   {سردسازی کنترل‌شده (کمینه‌سازی تنش پسماند)};

  \foreach \n/\nd in {1/a,2/b,3/c,4/d,5/e,6/f}{
    \node[num] at (\nd.west) {\n};
  }
  \foreach \i/\j in {a/b,b/c,c/d,d/e,e/f}{
    \draw[->, very thick, teal!60!black] (\i.south) -- (\j.north);
  }
\end{tikzpicture}}
\end{document}
